% ============================================================
% Chapter 25 --- Pullbacks, integration of forms and orientability
% Originally: content/week-12.tex in full (Corsin Week 12), together with
%             content/week-13.tex, \section{Submanifolds with boundary and induced
%             orientation} (Corsin Week 13), which belongs with orientability rather than
%             with the statement of Stokes' theorem it happened to follow.
% ============================================================
\chapter{Pullbacks, Integration of Forms and Orientability}
\label{ch:pullbacks}

% Transition: Claude Opus 5 (Medium)
To integrate a differential form over a submanifold one has to pull it back to a parameter
domain, and then check that the answer does not depend on which parametrization was chosen. It
does not, provided the parametrizations agree on a direction. That proviso is orientability, and
it is what this chapter is really about. The last section fixes the induced orientation on a boundary, which
is the convention that makes the sign in Stokes' theorem come out right.

% Originally: content/week-12.tex, \section{Pullback of forms} (Corsin Week 12)


\section{Pullback of forms}
\label{sec:pullback_of_forms}

% Source: Corsin Nick/Class Notes/Week 12.pdf, p. 5
\begin{definition}[Pullback of a form]
\label{def:pullback_of_forms}
Let $U \subseteq \mathbb{R}^m$, $V \subseteq \mathbb{R}^n$ be open, $F : U \to V$ a smooth
map, and $\omega \in \Omega^p(V)$ a differential $p$-form on $V$. The \newterm{pullback}
(\germanterm{Zur\"uckziehung}) of $\omega$ under $F$ is the differential $p$-form on $U$
defined by
\[(F^*\omega)_x(v_1,\ldots,v_p) := \omega_{F(x)}\big(DF_x(v_1),\ldots,DF_x(v_p)\big).\]
\end{definition}

\begin{ainote}
Corsin also draws a commutative diagram $TU \xrightarrow{dF} TV \xrightarrow{\omega} \mathbb{R}$
over $U \xrightarrow{F} V$, with tangent-bundle projections $\pi_U, \pi_V$, and labels it
\qt{not exam relevant}. Omitted here for that reason; the formula above is the operative
content.
\end{ainote}

The pullback satisfies two properties that make it the right tool for transporting forms
along smooth maps:
\begin{enumerate}[label=\textbf{(\arabic*)}]
  \item $F^*(\omega\wedge\theta) = F^*\omega \wedge F^*\theta$;
  \item $F^*(d\omega) = d(F^*\omega)$ (the pullback commutes with the exterior derivative).
\end{enumerate}

\begin{ainote}
This is precisely the operation used, informally, in the reparametrization-invariance
motivation for \cref{def:antisymmetric_k_linear_form}: writing $\phi^*\omega$ and
$(\phi\circ\Psi)^*\omega$ for the two sides of that calculation makes it read as $\int_D
\phi^*\omega = \int_V (\phi\circ\Psi)^*\omega = \int_V \Psi^*(\phi^*\omega)$, i.e.\ property (1)
applied to $\phi^*\omega$ and the fact that pullback along a composition is the composition of
pullbacks.
\end{ainote}

% Originally: content/week-12.tex, \section{Integration of differential forms on submanifolds} (Corsin Week 12)

\section{Integration of differential forms on submanifolds}
\label{sec:integration_of_forms}

% Source: Corsin Nick/Class Notes/Week 12.pdf, p. 6
If $\mathcal{M}$ is an $m$-dimensional submanifold of $\mathbb{R}^m$ --- i.e.\ simply an open
set --- and $\omega \in \Omega^m(\mathcal{M})$ an $m$-form, then
\[\int_{\mathcal{M}} \omega := \int_{\mathcal{M}} \omega_x(e_1,\ldots,e_m)\,dx;\]
in particular, for $f \in C^\infty(\mathcal{M})$,
\[\int_{\mathcal{M}} f\,dx^1\wedge\cdots\wedge dx^m = \int_{\mathcal{M}} f(x)\,dx,\]
recovering the ordinary Lebesgue integral of $f$.

\begin{definition}[Integral of a form over a parametrized submanifold]
\label{def:integral_of_form_over_submanifold}
If $L^k \subseteq \mathbb{R}^n$ is a $k$-dimensional submanifold, $n>k$, and $\phi : U \to L$
is a parametrization with $U \subseteq \mathbb{R}^k$ open, then the integral of $\omega \in
\Omega^k(L)$ is defined as
\[\int_L \omega := \int_U \phi^*\omega = \int_U \omega_{\phi(x)}\!\left(\frac{\partial\phi}
  {\partial x^1},\ldots,\frac{\partial\phi}{\partial x^k}\right) dx.\]
\end{definition}

In particular, for $f \in C^\infty(L)$ and $1\leq i_1<\cdots<i_k\leq n$,
\[\int_L f\,dx^{i_1}\wedge\cdots\wedge dx^{i_k} = \int_U f(x)\det\!\left(\frac{\partial
  \phi^{i_s}}{\partial x^t}\right)_{1\leq s,t\leq k} dx,\]
because, by property (4) of \cref{def:wedge_product},
\[dx^{i_1}\wedge\cdots\wedge dx^{i_k}(\partial_1\phi,\ldots,\partial_k\phi) = \det\!\left(
  \begin{smallmatrix} dx^{i_1}(\partial_1\phi) & \cdots & dx^{i_1}(\partial_k\phi) \\
  \vdots & \ddots & \vdots \\ dx^{i_k}(\partial_1\phi) & \cdots & dx^{i_k}(\partial_k\phi)
  \end{smallmatrix}\right) = \det\!\left(\frac{\partial\phi^{i_s}}{\partial x^t}\right)_{1\leq
  s,t\leq k}.\]

\begin{ainote}
Applying \cref{def:integral_of_form_over_submanifold} to $\omega_F^{n-1}$ or $\omega_F^2$ from
\cref{ch:differential_forms} recovers exactly the flux integrals of \cref{ch:flux} --- those
were, all along,
integrals of forms in this sense. And \cref{def:pullback_of_forms} is what makes the
motivating reparametrization-invariance calculation precise: $\int_L \omega$ does
not depend on the parametrization $\phi$ chosen, only on $L$ and its orientation.
\end{ainote}


% --- exercises relocated here from the dissolved sheets ---

% Originally: content/week-12.tex, \exercisesheet{12}, problem 12.9
% Source: exercises/Ex12_Analysis2_eng.pdf

\begin{exercise}[12.9 --- Change of variables with differential forms \textnormal{(important)}]
\label{ex:12.9}
Let $U,V \subseteq \mathbb{R}^n$ be open sets, and let $\psi = (\psi^1,\ldots,\psi^n) : U \to
V$ be a diffeomorphism with $\det\!\left(\frac{\partial\psi^j}{\partial x^i}\right) > 0$.
Define an $n$-differential form $\omega$ on $U$ by $\omega = d\psi^1\wedge d\psi^2\wedge\cdots
\wedge d\psi^n$.
\begin{enumerate}[label=\textbf{(\alph*)}]
  \item Show that $\omega = \det\!\left(\frac{\partial\psi^j}{\partial x^i}\right) dx^1\wedge
  \cdots\wedge dx^n$.
  \item From the relation in (a), rewrite the change of variables formula using differential
  forms.
\end{enumerate}
\end{exercise}

% Originally: content/week-12.tex, \section{Orientable manifolds} (Corsin Week 12)

\section{Orientable manifolds}
\label{sec:orientable_manifolds}

% Source: Corsin Nick/Class Notes/Week 12.pdf, p. 7
\begin{ainote}
The two definitions below are reproduced by Corsin directly from the lecture script (numbered
there as Definitions 14.68 and 14.65), rather than written out in his own hand; he adds the
margin annotations quoted and figured below.
\end{ainote}

\begin{definition}[Support of a $k$-form]
\label{def:support_of_k_form}
Let $U \subseteq \mathbb{R}^n$ be open and $\alpha \in \Omega^k(U)$. The \newterm{support}
(\germanterm{Tr\"ager}) of $\alpha$ is the closed set
\[\supp\alpha := \overline{\{x\in U : \alpha_x \neq 0\}} \subseteq \overline{U}.\]
We say $\alpha$ has \newterm{compact support} in $U$ if $\supp\alpha$ is bounded and contained
in $U$.
\end{definition}

\begin{definition}[Orientable manifold]
\label{def:orientable_manifold}
Let $M \subseteq \mathbb{R}^n$ be a $k$-dimensional $C^\infty$ submanifold. We say $M$ is
\newterm{orientable} (\germanterm{orientierbar}) if there exists a family of local
parametrizations $\phi_i : V_i \to M \cap U_i$, $i \in I$, whose images cover $M$, and such
that for every $i,j \in I$ with $M \cap U_i \cap U_j \neq \emptyset$, the transition map
$\widetilde\phi_i^{-1}\circ\widetilde\phi_j : \widetilde V_j \to \widetilde V_i$ has positive
Jacobian determinant:
\[\det D(\widetilde\phi_i^{-1}\circ\widetilde\phi_j) > 0 \qquad \text{on } \widetilde V_j.\]
Such a family $\{\phi_i\}$ is called an \newterm{atlas} of $M$; if it satisfies the
determinant condition above it is called an \newterm{orientation} of $M$. Once such a family
is chosen, $M$ is said to be \newterm{oriented}.
\end{definition}

\begin{ainote}
Corsin marks the family $\{\phi_i\}$ itself with the margin note \qt{is called an Atlas} next
to the covering condition, and the positivity condition with \qt{Atlas} again next to the
determinant inequality --- i.e.\ he is flagging that \qt{atlas} names the family regardless of
whether the determinant condition holds, while \qt{orientation} is the stronger notion once it
does. This matches the definition's own wording above.
\end{ainote}

Corsin illustrates the failure of the orientation condition with a hand-drawn counterexample:
two overlapping local parametrizations $\phi, \psi$ of the same patch of a surface, whose
coordinate frames $(\partial_1\phi,\partial_2\phi)$ and $(\partial_1\psi,\partial_2\psi)$ at a
shared point disagree in handedness.

% Sonnet 5 (Medium) --- FIG-W12-01
\begin{center}
\begin{tikzpicture}[>=Latex, scale=1.0]
  \draw[green!50!black, thick] (0,0) .. controls (1.5,0.9) and (3.5,0.6) .. (5,1.2)
    .. controls (5,2.4) and (3,2.6) .. (1.5,2.0) .. controls (0,1.6) and (0,0.8) .. (0,0);
  \draw[green!50!black, dotted] (0.4,1.0) .. controls (1.8,0.6) and (3.2,1.1) .. (4.4,1.6);
  \fill (2.1,1.3) circle (1.2pt);
  \draw[blue!70!black, thick, ->] (2.1,1.3) -- (3.3,1.55) node[right, font=\scriptsize]
    {$\partial_1\phi$};
  \draw[blue!70!black, thick, ->] (2.1,1.3) -- (1.7,2.15) node[above, font=\scriptsize]
    {$\partial_2\phi$};
  \draw[red, thick, ->] (2.1,1.3) -- (0.9,1.05) node[left, font=\scriptsize] {$\partial_1\psi$};
  \draw[red, thick, ->] (2.1,1.3) -- (2.55,0.35) node[below, font=\scriptsize] {$\partial_2\psi$};
  \node[font=\scriptsize] at (0.1,-0.35) {$\phi$};
  \node[font=\scriptsize] at (4.4,2.7) {$\psi$};
  \draw[blue!70!black, ->] (6.3,0.2) -- (7.3,0.2) node[right, font=\scriptsize] {$e_1$};
  \draw[blue!70!black, ->] (6.3,0.2) -- (6.3,1.2) node[above, font=\scriptsize] {$e_2$};
  \draw[red, ->] (8.6,0.2) -- (9.6,0.2) node[right, font=\scriptsize] {$e_1$};
  \draw[red, ->] (8.6,0.2) -- (8.6,1.2) node[above, font=\scriptsize] {$e_2$};
\end{tikzpicture}
\end{center}

\begin{ainote}
At the shared point, $(\partial_1\phi,\partial_2\phi)$ turns counterclockwise while
$(\partial_1\psi,\partial_2\psi)$ turns clockwise --- the two frames are mirror images of each
other. Corsin's caption: \qt{Here, $\det(D(\phi^{-1}\circ\psi)) < 0$!} This is exactly the atlas
$\{\phi,\psi\}$ failing \cref{def:orientable_manifold}: it covers the patch, but its transition
map has \emph{negative} Jacobian determinant there, so $\{\phi,\psi\}$ is an atlas but not an
orientation.
\end{ainote}

% Generator: Claude Sonnet 5 (Medium)
\begin{ainote}
\Cref{ex:12.3} and Exercise 12.6 (the latter not reproduced below, see the priority table)
invoke \newterm{Stokes' theorem} in its classical $\mathbb{R}^3$ form, which none of Corsin's
transcribed pages state explicitly. Recorded here for self-containedness, since the
exercise-sheet solutions below rely on it.
\end{ainote}

\begin{theorem}[Stokes' theorem, classical form in $\mathbb{R}^3$]
\label{thm:stokes_classical}
% Generator: Claude Sonnet 5 (Medium)
Let $S \subseteq \mathbb{R}^3$ be a compact, oriented $C^1$ surface with boundary $\partial S$
carrying the induced orientation, and let $F$ be a $C^1$ vector field on a neighborhood of
$S$. Then
\[\int_S \curl F \cdot d\vec{n} = \int_{\partial S} F \cdot d\vec{s}.\]
\end{theorem}

% Originally: content/week-13.tex, \section{Submanifolds with boundary and induced orientation} (Corsin Week 13)


\section{Submanifolds with boundary and induced orientation}
\label{sec:submanifolds_with_boundary}

% Source: Corsin Nick/Class Notes/Week 13.pdf, pp. 6--7
\begin{ainote}
Corsin recaps \cref{thm:local_parametrization} verbatim --- for every $p \in M$, an
open neighbourhood $V$ of $p$, an open set $U \subseteq \mathbb{R}^k$, and $\phi : U \to
V\cap M$ that is smooth, has $D\phi_x$ injective for all $x \in U$, and is a homeomorphism onto
$V\cap M$ --- before extending it to submanifolds with boundary below. Not repeated; see
\cref{sec:submanifolds}.
\end{ainote}

\begin{definition}[Submanifold with boundary]
\label{def:submanifold_with_boundary}
For a $k$-dimensional submanifold \emph{with boundary}, we require similarly that for a local
parametrization $\phi : U \to V\cap M$ as in \cref{thm:local_parametrization}, either $U
\subseteq \mathbb{R}^k$ is open, or
\[U = U' \cap \mathcal{H}^k, \qquad U' \subseteq \mathbb{R}^k \text{ open}, \qquad
  \mathcal{H}^k := \{(x_1,\ldots,x_k) \in \mathbb{R}^k : x_1 \leq 0\}\]
the (closed) \newterm{left half-space}. If $p \in \partial M$ and $\phi : U \subseteq
\mathcal{H}^k \to V\cap M$ is a local parametrization of $M$ around $p$, then $\phi|_{\partial
\mathcal{H}^k}$ (identifying $\partial\mathcal{H}^k = \{0\}\times\mathbb{R}^{k-1} \cong
\mathbb{R}^{k-1}$) is a local parametrization of $\partial M$ around $p$.
\end{definition}

% Sonnet 5 (Medium) --- FIG-W13-01
\begin{center}
\begin{tikzpicture}[>=Latex, scale=1.0]
  \draw[green!50!black, thick] (0,0) .. controls (0.3,1.3) and (1.3,2.0) .. (2.6,1.9)
    .. controls (3.9,1.8) and (4.6,1.0) .. (4.4,0.1) .. controls (3.5,-0.5) and (1.5,-0.6)
    .. (0,0);
  \fill (0.9,0.55) circle (1.1pt);
  \draw[red, dashed, thick] (-0.05,-0.55) .. controls (0.4,0.0) and (0.5,1.1) .. (0.15,1.55);
  \draw[red, thick, ->] (0.9,0.55) -- (1.75,0.75) node[right, font=\scriptsize] {$\partial_2\phi$};
  \draw[red, thick, ->] (0.9,0.55) -- (0.55,1.3) node[above, font=\scriptsize] {$\partial_1\phi$};
  \node[font=\scriptsize] at (0.6,-0.85) {$p \in \partial M$};
  \draw[->] (5.3,0.7) -- (6.3,0.7) node[right, font=\scriptsize] {};
  \node[font=\scriptsize] at (5.8,1.0) {$\phi$};
  \draw[red, thick] (7.1,0.1) rectangle (8.4,1.4);
  \draw[red, thick, dashed] (7.1,0.1) -- (7.1,1.4);
  \draw[red, ->] (7.1,0.75) -- (7.7,0.75) node[right, font=\scriptsize] {$e_1$};
  \draw[red, ->] (7.1,0.75) -- (7.1,1.25) node[above, font=\scriptsize] {$e_2$};
  \node[font=\scriptsize] at (7.1,-0.15) {$0$};
  \node[font=\scriptsize] at (7.75,-0.2) {$U \subseteq \mathcal{H}^2$};
\end{tikzpicture}
\end{center}

\begin{definition}[Induced orientation]
\label{def:induced_orientation}
If $\{\partial_1\phi,\ldots,\partial_k\phi\}$ is the oriented basis of $T_pM$ given by a local
parametrization $\phi$ around $p \in \partial M$, then $\{\partial_2\phi,\ldots,\partial_k
\phi\}$ is the \newterm{induced} oriented basis of $T_p(\partial M)$ around $p$.
\end{definition}

\subsection{Example: cylinder}
\label{sec:cylinder_orientation_example}

% Source: Corsin Nick/Class Notes/Week 13.pdf, pp. 8--9
Consider the cylindrical surface
\[C = \{(x,y,z) \in \mathbb{R}^3 : y^2+z^2 = r^2,\ 0\leq x\leq L\},\]
a $2$-dimensional submanifold with boundary, and the two parametrizations $\phi,\psi :
[-L,0]\times[0,2\pi] \to C$ (both domains $\subseteq \mathcal{H}^2$),
\[\phi(x,\theta) = \begin{pmatrix} -x \\ r\cos\theta \\ -r\sin\theta \end{pmatrix}, \qquad
  \psi(x,\theta) = \begin{pmatrix} x+L \\ r\cos\theta \\ r\sin\theta \end{pmatrix}.\]
Both cover all of $C$: as the local $x$-coordinate ranges over $[-L,0]$, $-x$ and $x+L$ both
range over $[0,L]$.

The two are positively oriented with respect to each other: writing out the composition,
\[\phi^{-1}\circ\psi(x,\theta) = \phi^{-1}(x+L,\,r\cos\theta,\,r\sin\theta) = (-x-L,-\theta)
  = -(x+L,\theta),\]
so $D(\phi^{-1}\circ\psi) = \begin{pmatrix}-1&0\\0&-1\end{pmatrix}$ and $\det D(\phi^{-1}
\circ\psi) = 1 > 0$: $\{\phi,\psi\}$ is a valid orientation of $C$ by
\cref{def:orientable_manifold}.

But restricted to the line $\partial\mathcal{H}^2 = \{(x,\theta) : x=0\}$, $\phi$ and $\psi$
parametrize \emph{opposite} boundary components of the cylinder --- $\phi(0,\cdot)$ traces the
circle at $x=0$, $\psi(0,\cdot)$ the circle at $x=L$ --- and \cref{def:induced_orientation}
gives them opposite rotational senses when both are drawn in the same $(y,z)$-view:

% Sonnet 5 (Medium) --- FIG-W13-02
\begin{center}
\begin{tikzpicture}[>=Latex, scale=0.95]
  \draw[blue!70!black, thick] (0,2.4) rectangle (4.5,3.5);
  \draw[blue!70!black, thick] (0,2.4) circle (0.55 and 0.55);
  \draw[blue!70!black, thick] (4.5,2.95) ellipse (0.35 and 0.55);
  \draw[red, thick, ->] (-0.5,3.5) arc (150:490:0.55);
  \node[red, font=\scriptsize] at (-0.85,3.5) {$t\mapsto\phi(0,t)$};
  \node[font=\scriptsize] at (0,1.7) {\textbf{left end} ($x=0$)};

  \draw[blue!70!black, thick] (6.5,2.4) rectangle (11,3.5);
  \draw[blue!70!black, thick, dashed] (6.5,2.95) ellipse (0.35 and 0.55);
  \draw[blue!70!black, thick] (11,2.4) circle (0.55 and 0.55);
  \draw[orange, thick, ->] (11.5,2.4) arc (-30:330:0.55);
  \node[orange, font=\scriptsize] at (11.9,3.5) {$t\mapsto\psi(0,t)$};
  \node[font=\scriptsize] at (8.75,1.7) {\textbf{right end} ($x=L$)};
\end{tikzpicture}
\end{center}

\begin{ainote}
This is not a contradiction, and Corsin does not present it as one: $\{\phi,\psi\}$ is a
perfectly good orientation of $C$ as a whole (the overlap check above is what
\cref{def:orientable_manifold} demands, and it passes). What the picture shows is that a
\emph{consistently} oriented cylinder induces boundary orientations on its two ends that look
like opposite rotational senses when both circles are drawn from the same external viewpoint
--- exactly the familiar fact, from the divergence theorem, that the outward normal at the two
ends of a solid cylinder points in opposite $x$-directions. Induced orientation is doing its
job correctly here, not failing.
\end{ainote}

% Originally: content/week-12.tex, \section{Solutions}

\section{Solutions}
\label{sec:pullbacks_solutions}

\begin{exercisesolution}[Solution of \cref{ex:12.9}]
\textbf{(a)} Writing each $d\psi^j$ in the coordinate basis, $d\psi^j = \sum_{i=1}^n
\frac{\partial\psi^j}{\partial x^i}\,dx^i$, so
\[\omega = \left(\sum_{i_1=1}^n \frac{\partial\psi^1}{\partial x^{i_1}}dx^{i_1}\right)\wedge
  \cdots\wedge\left(\sum_{i_n=1}^n \frac{\partial\psi^n}{\partial x^{i_n}}dx^{i_n}\right)
  = \sum_{i_1,\ldots,i_n} \frac{\partial\psi^1}{\partial x^{i_1}}\cdots
  \frac{\partial\psi^n}{\partial x^{i_n}}\,dx^{i_1}\wedge\cdots\wedge dx^{i_n}\]
by bilinearity of $\wedge$ (property (1) of \cref{def:wedge_product}). Any term with a
repeated index vanishes, since $dx^i\wedge dx^i=0$ by antisymmetry, so only tuples
$(i_1,\ldots,i_n) = (\sigma(1),\ldots,\sigma(n))$ for a permutation $\sigma \in S_n$
contribute; using $dx^{\sigma(1)}\wedge\cdots\wedge dx^{\sigma(n)} = \sgn(\sigma)\,dx^1\wedge
\cdots\wedge dx^n$,
\[\omega = \left(\sum_{\sigma\in S_n} \sgn(\sigma)\,\frac{\partial\psi^1}{\partial
  x^{\sigma(1)}}\cdots\frac{\partial\psi^n}{\partial x^{\sigma(n)}}\right) dx^1\wedge\cdots
  \wedge dx^n.\]
The coefficient is exactly the Leibniz expansion of $\det(D\psi)$, so
\[\boxed{\ d\psi^1\wedge\cdots\wedge d\psi^n = \det(D\psi)\,dx^1\wedge\cdots\wedge dx^n.\ }\]

\textbf{(b)} Let $\eta = f(y)\,dy^1\wedge\cdots\wedge dy^n$ be an $n$-form on $V$. Its
pullback under $\psi$ is $\psi^*\eta = f(\psi(x))\,d\psi^1\wedge\cdots\wedge d\psi^n$, which by
part (a) equals $f(\psi(x))\det(D\psi(x))\,dx^1\wedge\cdots\wedge dx^n$. By
\cref{def:integral_of_form_over_submanifold} (applied with $U$ itself as the $n$-dimensional
submanifold),
\[\int_U \psi^*\eta = \int_U f(\psi(x))\det(D\psi(x))\,dx.\]
Since $\det(D\psi) = \lvert\det(D\psi)\rvert>0$ by hypothesis, the right-hand side is exactly
the ordinary change-of-variables formula, so
\[\int_U \psi^*\eta = \int_U f(\psi(x))\,\lvert\det(D\psi(x))\rvert\,dx = \int_V f(y)\,dy
  = \int_V \eta.\]
Thus $\boxed{\int_U \psi^*\eta = \int_V \eta}$: the change-of-variables formula is exactly the
statement that integration of forms is invariant under pullback by an orientation-preserving
diffeomorphism.
\end{exercisesolution}

% Verified: solved independently, both parts match Sol12, including pullback-invariance.

\begin{exercisesolution}[Solution to \cref{ex:ai_pullback_area_form}]
% Generator: Claude Opus 5 (High)
\textbf{(a)} The pullback of a coordinate differential is the differential of the corresponding
component, so from $x = r\cos\theta$ and $y = r\sin\theta$,
\[\Phi^*(dx) = \cos\theta\,dr - r\sin\theta\,d\theta,
  \qquad
  \Phi^*(dy) = \sin\theta\,dr + r\cos\theta\,d\theta.\]
Wedging, and using $dr\wedge dr = d\theta\wedge d\theta = 0$ together with
$d\theta\wedge dr = -\,dr\wedge d\theta$,
\[\Phi^*(dx\wedge dy)
  = \big(\cos\theta\cdot r\cos\theta\big)\,dr\wedge d\theta
    + \big(-r\sin\theta\cdot\sin\theta\big)\,d\theta\wedge dr
  = r\big(\cos^2\theta + \sin^2\theta\big)\,dr\wedge d\theta
  = r\,dr\wedge d\theta.\]
The factor is exactly the $r$ tabulated for polar coordinates in
\cref{rem:coordinate_jacobians_table}, obtained here without computing a determinant: the
antisymmetry of the wedge did the cofactor expansion.

\textbf{(b)} They measure different things, and the absolute value is the difference.
\Cref{thm:change_of_variables} is a statement about \emph{volume}, which is unsigned: a region
and its mirror image have the same area, so only $\lvert\det\Jac\Phi\rvert$ can appear. The
pullback of a top-degree form is a statement about \emph{oriented} volume, which carries a sign,
and it records $\det\Jac\Phi$ itself.

The swap $\Phi(u_1,u_2) = (u_2,u_1)$ separates them. Its Jacobian is
$\left(\begin{smallmatrix}0&1\\1&0\end{smallmatrix}\right)$ with determinant $-1$, so
\[\Phi^*(dx\wedge dy) = du_2\wedge du_1 = -\,du_1\wedge du_2,\]
while $\lvert\det\Jac\Phi\rvert = 1$ and the change of variables formula does not change at all.
The form notices the reversal; the volume cannot.

\textbf{(c)} \Cref{ex:12.9} concludes that $\int_U\psi^*\eta = \int_V\eta$, with no sign. Drop
the hypothesis $\det > 0$ and take $\psi$ to be the swap above: the pullback acquires a factor
$-1$, so the two integrals differ by a sign and the conclusion is false as stated. The honest
statement without the hypothesis carries $\sgn(\det\Jac\psi)$.

\Cref{thm:change_of_variables} is untroubled by the same map because it never claimed anything
about orientation. This is the precise sense in which the positivity hypothesis in
\cref{ex:12.9} is not a technical convenience: it is the price of upgrading an unsigned identity
to a signed one.
\end{exercisesolution}

\begin{exercisesolution}[Solution to \cref{ex:ai_area_form_over_hemisphere}]
% Generator: Claude Opus 5 (High)
\textbf{(a)} Along $\phi$ the first two coordinates are the parameters themselves, $x = u$ and
$y = v$, so $\phi^*(dx) = du$ and $\phi^*(dy) = dv$, giving
\[\phi^*(dx\wedge dy) = du\wedge dv.\]
The third coordinate never enters, which is the whole reason this integral is easy. It remains to
check that $\phi$ induces the outward orientation: for a graph $z = g(u,v)$,
\[\partial_u\phi\times\partial_v\phi = (-g_u,\ -g_v,\ 1),\]
whose third component is $+1$, so the induced normal points \emph{upward}, and on the upper
hemisphere upward is outward. Hence
\[\int_{S^+} dx\wedge dy = \int_D du\wedge dv = \vol_2(D) = \pi.\]

\textbf{(b)} The pullback computation is identical, since $\psi$ also has $x = u$ and $y = v$,
so again $\psi^*(dx\wedge dy) = du\wedge dv$. But $\psi$ is still a graph, so its induced normal
still points upward, and on the \emph{lower} hemisphere upward is \emph{inward}. So $\psi$ is
orientation-reversing for the prescribed outward orientation, and
\cref{def:integral_of_form_over_submanifold} requires a parametrization compatible with the
chosen orientation. Correcting by the sign,
\[\int_{S^-} dx\wedge dy = -\int_D du\wedge dv = -\pi.\]

\textbf{(c)} The equator is a null set, so the two halves add:
\[\int_{S^2} dx\wedge dy = \pi + (-\pi) = 0.\]
For the second route, $d(x\,dy) = dx\wedge dy$, and $S^2$ is a compact orientable $2$-manifold
with $\partial S^2 = \emptyset$, so \cref{thm:stokes_general} gives
\[\int_{S^2} dx\wedge dy = \int_{S^2} d(x\,dy) = \int_{\partial S^2} x\,dy = 0.\]

\begin{remark}[$dx\wedge dy$ measures the shadow]
\label{rem:area_form_measures_the_shadow}
Part \textbf{(a)} says something worth keeping: integrating $dx\wedge dy$ over an oriented
surface in $\mathbb{R}^3$ returns the \emph{signed} area of its shadow on the $xy$-plane. The
hemisphere has area $2\pi$, but its shadow is the unit disc, and $\pi$ is what the computation
returns. A $2$-form does not measure surface area; the Gram determinant of
\cref{ch:gram_determinant} does that, and the two answers differ by exactly the tilt of the
surface.

Read that way, part \textbf{(c)} is obvious before any computation. The sphere's two halves cast
the same shadow, and a closed surface presents each patch of shadow once from above and once from
below, so everything cancels. That is the geometric content of $\int_M d\omega = 0$ on a closed
$M$, and \cref{q:quiz_q14} draws the other standard consequence from it.
\end{remark}
\end{exercisesolution}

